\section{Gabber's Involutivity Theorem}
Let $B$ be a commutative $\C$-algebra.

\begin{definition}
  \leavevmode
  \begin{enumerate}
    \item A $\C$-bilinear map $\{\,,\,\}\colon B\otimes_{\C}B\to B$ is a \textbf{Poisson bracket}
          if it is a Lie bracket and satisfies the Leibniz rule $\{a_1a_2,b\}=a_1\{a_2,b\}+a_2\{a_1,b\}$ for all $a_1,a_2,b\in B$.
    \item If $B$ is Poisson, a closed subset $V\subseteq\Spec B$ with radical ideal $I_V$ is
          \textbf{coisotropic} (or \textbf{involutive}) if $\{I_V,I_V\}\subseteq I_V$.
  \end{enumerate}
\end{definition}

\noindent $\Gamma_\bullet A : \text{a } \mathbb{Z}_{\ge 0}\text{-filtered } \mathbb{C}\text{-algebra}$
\begin{definition}
  \leavevmode
  \begin{enumerate}
    \item $\Gamma_\bullet A$ is called \textbf{almost commutative} if $\operatorname{gr}^\Gamma_\bullet A$ is a commutative $\mathbb{C}$-algebra of finite type.
    \item For $x \in \Gamma_k A$, we use $\sigma_k(x) \text{ to denote its image in } \frac{\Gamma_k A}{\Gamma_{k-1} A}$.
    \item For $x \in \Gamma_k A$, $y \in \Gamma_{\ell} A$, we define $\{ \sigma_k(x), \sigma_{\ell}(y) \} = \sigma_{k+\ell-1}(xy - yx)$
          (Since $\operatorname{gr}_\bullet A$ is commutative, $xy - yx \in \Gamma_{k+\ell-1} A$)
  \end{enumerate}
\end{definition}

\begin{lemma}[Lemma~6.3]
  $(\gr^{\Gamma}_\bullet A,\{\,\,\})$ is a Poisson algebra.
\end{lemma}
\begin{proof}
  Homework.
\end{proof}

\begin{example}
  For a smooth affine $X$, the identification $\gr \D_X\cong \Sym^{\bullet}\T_X$ transports the
  Poisson structure on $\gr \D_X$ to $T^{\ast}X$.
\end{example}

If $\mathcal M$ is a finitely generated $A$-module with a good filtration, the characteristic
variety $\Ch(\mathcal M):=\Supp(\gr^{\Gamma}_\bullet\!\mathcal M)$ is a closed subset of
$\Spec\gr^{\Gamma}_\bullet A$.

\begin{theorem}[Gabber's~Involutivity~I]\label{thm:gabber-I}
  If $\mathcal M$ is a finitely generated $A$-module, then $\Ch(\mathcal M)\subseteq\Spec\gr^{\Gamma}_\bullet A$ is involutive.
\end{theorem}

\begin{corollary}
  For a finitely generated $\mathcal{D}_X$-module $\mathcal{M}$ on a smooth affine variety $X$, the characteristic variety $\operatorname{Ch}(\mathcal{M})$ inside $T^{\ast}X$ is conic and involutive. In particular $\dim \operatorname{Ch}(\mathcal{M}) \ge n$.
\end{corollary}

To prove Theorem~\ref{thm:gabber-I}, we introduce Gabber rings.

\begin{definition}
  Let $\mathbb E:=\C[u]/(u^2)$ be the \textbf{ring of dual numbers}. An $\mathbb E$-algebra $T$ (not
  necessarily commutative) is a \textbf{Gabber ring} if
  \begin{enumerate}
    \item $T$ is flat over $\mathbb E$, and
    \item $\overline T:=T/uT$ is a commutative $\C$-algebra of finite type.
  \end{enumerate}
\end{definition}

\begin{lemma}\label{E-flatness}
  For a $T$-module $\underline{\mathcal M}$, the following are equivalent:
  \begin{enumerate}
    \item $\underline{\mathcal M}$ is $\mathbb E$-flat.
    \item The map $\underline{\mathcal M}/u\underline{\mathcal M} \to u\underline{\mathcal M},\ \bar{m} \mapsto um$ is an isomorphism.
  \end{enumerate}
\end{lemma}
\begin{proof}
  Recall that for an $R$-module $N$ over a commutative ring $R$,
  \[
    N \text{ is } R\text{-flat} \iff \forall \text{ proper ideal } I \subseteq R, \quad I \otimes_R N \to N \text{ is injective.}
  \]

  So $\underline{M}$ is $\mathbb{E}$-flat $\iff (u) \otimes_{\mathbb{E}} \underline{M} \hookrightarrow \underline{M}$
  (since $(u)$ is the only non-trivial ideal of $\mathbb{E}$).

  Note that the image of the map is $(u)\underline{M} = u \cdot \underline{M}$.
  Consider the map $\mathbb{E}/(u) \to (u)$ given by $1 \mapsto u$. This map is an $\mathbb{E}$-module isomorphism.

  Thus, we have an $\mathbb{E}$-module isomorphism:
  \[
    \underline{M}/(u)\underline{M} = \mathbb{E}/(u) \otimes_{\mathbb{E}} \underline{M} \xrightarrow{\sim} (u) \otimes_{\mathbb{E}} \underline{M}.
  \]

  Therefore, $\underline{M}$ is $\mathbb{E}$-flat $\iff$ the composite map $\underline{M}/(u)\underline{M} \longrightarrow u \cdot \underline{M} $
  is an isomorphism.
\end{proof}

Whenever $\underline{\mathcal M}$ is $\mathbb E$-flat and $\overline a,\overline b\in\overline T$, we can pick
lifts $a,b\in T$ and define $\{\overline a,\overline b\} := \overline c$ where $c$ solves $u\cdot c = ab-ba$. Since $\overline{T}$ is commutative, such a $c$ exists.

\begin{lemma}\label{lem:poisson}
  $(\overline T,\{\,\,\})$ is a Poisson algebra.
\end{lemma}
\begin{proof}
  Skew-symmetry is immediate from the definition. Take lifts $a,b,c\in T$ of
  $\overline a,\overline b,\overline c$. Direct calculation gives
  \begin{align*}
    \{\overline a,\{\overline b,\overline c\}\}
     & = \overline{a(bc-cb) - (bc-cb) a}, \\
    \{\overline b,\{\overline c,\overline a\}\}
     & = \overline{b(ca-ac) - (ca-ac) b}, \\
    \{\overline c,\{\overline a,\overline b\}\}
     & = \overline{c(ab-ba) - (ab-ba) c}.
  \end{align*}
  Adding the three lines shows Jacobi. For the Leibniz rule we compute
  \begin{align*}
    \{\overline a\,\overline a',\overline b\}
     & = \overline{aa' b - ba a'}
    = \overline{a(a' b - b a')} + \overline{(ab - ba)a'}                                        \\
     & = \overline a\,\{\overline a',\overline b\} + \{\overline a,\overline b\}\,\overline a',
  \end{align*}
  and similarly in the second argument. Hence $(\overline T,\{\,\,\})$ is Poisson.
\end{proof}

\begin{theorem}[Gabber's~Involutivity~II]\label{thm:gabber-II}
  Let $T$ be a Gabber ring and $\underline{\mathcal M}$ a finitely generated $\mathbb E$-flat $T$-module. If
  $I:=\rad\Ann_{\overline T}(\overline{\underline{\mathcal M}})$, then $\{I,I\}\subseteq I$.
\end{theorem}

\begin{proof}[Theorem~\ref{thm:gabber-II} implies Theorem~\ref{thm:gabber-I}]
  Put $\Rees^{\Gamma}_\bullet(A):=\bigoplus_{p\ge0}\Gamma_pA\,u^p\subseteq A[u]$.  This is a finitely generated
  $\C[u]$-algebra, and by construction it is torsion free over $\C[u]$; since $\C[u]$ is a PID, the
  module $\Rees^{\Gamma}_\bullet(A)$ is therefore flat over $\C[u]$.  Write
  \[
    T := \Rees^{\Gamma}_\bullet(A)\otimes_{\C[u]}\mathbb E
    = \frac{\Rees^{\Gamma}_\bullet(A)}{u^2\Rees^{\Gamma}_\bullet(A)},
    \qquad
    \overline T \cong \frac{\Rees^{\Gamma}_\bullet(A)}{u\,\Rees^{\Gamma}_\bullet(A)}
    = \gr^{\Gamma}_\bullet A.
  \]
  Thus $T$ is a Gabber ring with reduction $\overline T=\gr^{\Gamma}_\bullet A$.  If $F_\bullet\underline{\mathcal M}$ is a
  good filtration on a finitely generated $A$-module, its Rees module
  $\Rees^{F}_\bullet(\underline{\mathcal M})=\bigoplus_{p\ge0}F_p\underline{\mathcal M}\,u^p$ is again torsion free over $\C[u]$,
  hence flat.  After base change we obtain the $\mathbb E$-flat $T$-module
  \[
    \underline{\mathcal{M}} := \Rees^{F}_\bullet(\underline{\mathcal M})\otimes_{\C[u]}\mathbb E
    = \frac{\Rees^{F}_\bullet(\underline{\mathcal M})}{u^2\Rees^{F}_\bullet(\underline{\mathcal M})},
  \]
  whose reduction modulo $u$ is
  $\overline{\underline{\mathcal{M}}}=\Rees^{F}_\bullet(\underline{\mathcal M})/u\,\Rees^{F}_\bullet(\underline{\mathcal M})\cong\gr^{F}_\bullet\underline{\mathcal M}$.  Now
  Theorem~\ref{thm:gabber-II} applied to $(T,\underline{\mathcal{M}})$ says that
  $\Supp(\gr^{F}_\bullet\underline{\mathcal M})\subseteq\Spec\gr^{\Gamma}_\bullet A$ is involutive.

  \begin{lemma}\label{lem:rees}
    For $\overline T=\gr^{\Gamma}_\bullet A$ and $T=\frac{\Rees^{\Gamma}_\bullet(A)}{u^2\Rees^{\Gamma}_\bullet(A)}$
    (the $\mathbb E$-algebra that appeared above), the induced bracket coincides with the intrinsic
    bracket defined earlier on $\gr^{\Gamma}_\bullet A$.
  \end{lemma}
  \begin{proof}
    The editor of this lecture note think it is straightforward. And it is assigned as homework.
  \end{proof}

  Then we complete the deduction of Theorem~\ref{thm:gabber-I} from Theorem~\ref{thm:gabber-II}.
\end{proof}

\begin{proof}[Proof of Theorem~\ref{thm:gabber-II}]



  Write
  \[
    I=\rad\Ann_{\overline T}(\overline{\underline{\mathcal M}})
    = \bigcap_{\substack{
        \text{minimal prime ideal } \\
        \mathfrak p_i \supseteq \Ann_{\overline T}(\overline{\underline{\mathcal M}})
      }} \mathfrak p_i.
  \]
  Since $\{I,I\}\subseteq I$ iff $\{\mathfrak p_i,\mathfrak p_i\}\subseteq \mathfrak p_i$ for every minimal prime, we fix one such $\mathfrak p$ and work locally.

  \begin{lemma}\label{lem:noether}
    Let $R$ be a finite type $k$-algebra of dimension $\ell+d$ and $\mathfrak p\subset R$ a prime of codimension
    $\ell$. There exist algebraically independent elements $x_1,\dots,x_\ell,y_1,\dots,y_d\in R$ such that
    \begin{enumerate}
      \item $R$ is finite over $k[x_1,\dots,x_\ell,y_1,\dots,y_d]$.
      \item $\mathfrak p\cap k[x_1,\dots,x_\ell,y_1,\dots,y_d]=(x_1,\dots,x_\ell)$.
      \item There exists $0\ne f\in k[y_1,\dots,y_d]$ for which $(R/\mathfrak p)_f$ is free over $k[y_1,\dots,y_d]_f$.
    \end{enumerate}
    \noindent
    This is a version of the Noether Normalization Theorem (see, for instance, Eisenbud,
    \textbf{Commutative Algebra with a View Toward Algebraic Geometry}).
  \end{lemma}

  \begin{proof}[Sketch of Lemma~\ref{lem:noether}]
    By \cite[Theorem 13.3(Noether Normalization).]{eisenbud1995commutative} We proved (1) and (2).

    To prove (3), we use the following result called \textbf{Generic Freeness},
    This notion will be used again in the proof of Theorem~\ref{thm:gabber-II}:

    \begin{lemma}[{\cite[\href{https://stacks.math.columbia.edu/tag/051R}{Tag 051R}]{stacks-project} or \cite[Theorem 14.4(Grothendieck's Generic Freeness Lemma).]{eisenbud1995commutative}}]\label{algebra-lemma-generic-flatness-Noetherian}
      \mbox{}\\
      Let $R \to S$ be a ring map.
      Let $M$ be an $S$-module.
      Assume
      \begin{enumerate}
        \item $R$ is Noetherian,
        \item $R$ is a domain,
        \item $R \to S$ is of finite type, and
        \item $M$ is a finite type $S$-module.
      \end{enumerate}
      Then there exists a nonzero $f \in R$ such that
      $M_f$ is a free $R_f$-module.% Lemma content goes here
    \end{lemma}

    Using this Lemma~\ref{algebra-lemma-generic-flatness-Noetherian}, we proved (3) of Lemma~\ref{lem:noether}.
  \end{proof}



  Choose a minimal prime $\mathfrak p$ of $\overline T$ containing $\Ann_{\overline T}(\overline{\underline{\mathcal M}})$
  and write $d=\dim \mathfrak p$. Since $\overline T$ is a finite type $\C$-algebra, we have primary decomposition that is:
  \[
    \rad\Ann_{\overline T}(\overline{\underline{\mathcal M}})
    = \mathfrak p \cap \Bigl(\bigcap_{i\ne j} \mathfrak p_i\Bigr).
  \]
  and we assume $\mathfrak p$ and $\mathfrak p_i$ are all distinct minimal primes.

  Choose $g\in\bigcap_{i\ne j}\mathfrak p_i\smallsetminus \mathfrak p$. Localizing at $g$ removes all
  minimal primes except $\mathfrak p$, so $(\overline{\underline{\mathcal M}})_{g}$ is killed by some power
  $\mathfrak p_{g}^\ell$.

  Since $\overline{\underline{\mathcal M}}$ is finitely generated over $\overline T$, by repeatedly applying Lemma~\ref{algebra-lemma-generic-flatness-Noetherian}, we get:
  \[
    \exists g,\, \forall k,\, \frac{\mathfrak p^k_g \overline{\underline{\mathcal M}}}{\mathfrak p^{k+1}_g \overline{\underline{\mathcal M}}} \text{ is free over } \overline{T}_g/\mathfrak p_g.
  \]
  (One should note that $\overline{T}/ \mathfrak p$ is a noetherian domain, and $\mathfrak p^k \overline{\underline{\mathcal M}}/ \mathfrak p^{k+1} \overline{\underline{\mathcal M}}$ is finitely generated over $\overline{T}/ \mathfrak p$.)


  Apply Lemma~\ref{lem:noether} to the finite-type algebra $R=\overline T_{g}$. We obtain algebraically
  independent elements $y_1=t_1,\dots,y_d=t_d\in\overline T_{g}$ and some $h\in\C[t_1,\dots,t_d]$ such that both $\overline T_{g}$ and $(\overline T_{g}/\mathfrak p_{g})_h$ are finite and free
  over $S_h:=\C[t_1,\dots,t_d]_h$ (i.e. $S:=\C[t_1,\dots,t_d]$). Replacing $\bar f$ by $g h$ if necessary (and keeping the notation $\bar f$
  for the product).

  With this choice we have:
  $\mathfrak p_{\bar f}^\ell\overline{\underline{\mathcal M}}_{\bar f}=0$ for some $\ell>0$, and each successive
  quotient $\mathfrak p_{\bar f}^i\overline{\underline{\mathcal M}}_{\bar f} / \mathfrak p_{\bar f}^{i+1}\overline{\underline{\mathcal M}}_{\bar f}$
  is a free $S_h$-module. \\

  For later use, fix $S_h$-bases
  $\{\overline m_{i,j}\}$ such that for each $i$, $\{\overline m_{i,j}\}$(fix $i$, and change $j$) is a $S_h$-basis of
  $\mathfrak p_{\bar f}^i\overline{\underline{\mathcal M}}_{\bar f}$.
  This basis $\{\overline m_{i,j}\}$ can also be regarded as a $S_h$-basis of $\overline{\underline{\mathcal M}}_{\bar f}$
  (The $S_h$-action on $\overline{\underline{\mathcal M}}_{\bar f}$ is through $\overline{T}_{\bar f}$, not through $\overline{T}_{\bar f}/\mathfrak p_{\bar f}$).\\

  Relabelling we obtain:
  \[
    \underbrace{\overline m_{k_0},\dots,\overline m_{k_1-1}}_{\text{basis of }\overline{\underline{\mathcal M}}_{\bar f}/\mathfrak p_{\bar f}\overline{\underline{\mathcal M}}_{\bar f}},\quad
    \underbrace{\overline m_{k_1},\dots,\overline m_{k_2-1}}_{\text{basis of }\mathfrak p_{\bar f}\overline{\underline{\mathcal M}}_{\bar f}/\mathfrak p_{\bar f}^2\overline{\underline{\mathcal M}}_{\bar f}},\quad
    \ldots,\quad
    \underbrace{\overline m_{k_{\ell-1}},\dots,\overline m_{k_\ell-1}}_{\text{basis of }\mathfrak p_{\bar f}^{\ell-1}\overline{\underline{\mathcal M}}_{\bar f}}.
  \]

  We then take a lift $f\in T$ of $\bar f\in \overline T$.

  \begin{lemma}[Lemma~6.12]\label{lem:matrix}
    Let $a,b\in T$ have images $\bar a,\bar b\in \mathfrak p$. Then there exist integers
    $n_1,n_2,n_3\ge0$ and elements $e_{ij}\in T$ with $\overline{e_{ij}} \in S\subseteq \overline{T}$ such that
    \begin{enumerate}
      \item $f^{n_1}[f^{n_2}a,f^{n_3}b] m_i = u\sum_j e_{ij} m_j$ for every $i$, and
      \item $\sum_j \overline{e_{ij}} = 0$ for every $i$.
    \end{enumerate}
  \end{lemma}
  \begin{proof}
    If $k_{j_1} \le i < k_{j_1+1}$ for some $j_1$, then we have
    \[
      f^{n_2} a m_i = \sum_{j \ge k_{j_1+1}} u_{ij} m_j + u \sum_j \alpha_{ij} m_j,\quad \bar{u}_{ij}, \bar{\alpha}_{ij} \in S
    \]


    \begin{remark}
      By abuse of notation, let us clarify how to get this expression, this equality makes sense in $\underline{\mathcal{M}}$:

      \begin{enumerate}
        \item \textbf{Filtration:}
              $\overline{m_i} \in \mathfrak{p}^{k_i}_{\bar f}\overline{\underline{\mathcal{M}}}_{\bar f}/\mathfrak{p}^{k_i+1}_{\bar f}\overline{\underline{\mathcal{M}}}_{\bar f}$. Then $\overline{a}\overline{m_i} \in \mathfrak{p}^{k_i+1}_{\bar f}\overline{\underline{\mathcal{M}}}_{\bar f}$. So we have $\overline{a}\overline{m_i} = \sum_{j \ge k_{j_1+1}} \overline{u_{ij}}\overline{m_j} \in \overline{\underline{\mathcal{M}}}_{\bar f}$.

        \item \textbf{Localization:}
              So we have $\overline{f^{n_2} a m_i} = \sum_{j \ge k_{j_1+1}} \overline{u_{ij}} \overline{m_j} \in \overline{\underline{\mathcal{M}}}$, here we canceled the denominators of $\overline{u_{ij}}$, and use the same symbol to denote the new $\overline{u_{ij}}$.

        \item \textbf{Flatness:}
              By Lemma~\ref{E-flatness}, we have the isomorphism $\overline{\underline{\mathcal{M}}} \xrightarrow{\sim} u \underline{\mathcal{M}}$. So every element in $u\underline{\mathcal{M}}$ can be expressed by the form
              \[
                u \cdot \text{ elements lifted from }\overline{\underline{\mathcal{M}}}.
              \]
              So we have $f^{n_2} a m_i = \sum_{j \ge k_{j_1+1}} u_{ij} m_j + u \sum_j \alpha_{ij} m_j\quad  \in \underline{\mathcal{M}}$.
      \end{enumerate}
    \end{remark}



    Similarly,
    \[
      f^{n_3} b m_i = \sum_{j \ge k_{j_1+1}} v_{ij} m_j + u \sum_j \beta_{ij} m_j,\quad \bar{v}_{ij}, \bar{\beta}_{ij} \in S
    \]
    Then, since $u^2=0$,
    \begin{align*}
      f^{n_2} a f^{n_3} b m_i = \sum_{j \ge k_{j_1+1}} \left( v_{ij} f^{n_2} a + [f^{n_2}a, v_{ij}] \right) m_j + u \sum_j \beta_{ij} f^{n_2} a m_j
    \end{align*}

    \noindent Take $U = [u_{ij}]$, $V = [v_{ij}]$, $\alpha = [\alpha_{ij}]$, $\beta = [\beta_{ij}]$, we get:

    \begin{align*}
      f^{n_2} a f^{n_3} b m_i
      = \sum_{j \ge k_{j_1+1}} \left( (VU)_{ij} + [f^{n_2}a, v_{ij}] \right) m_j
      + u \sum_j \left( (V \cdot \alpha)_{ij} + (\beta \cdot U)_{ij} \right) m_j
    \end{align*}



    \noindent Similarly,
    \begin{align*}
      f^{n_3} b f^{n_2} a m_i
      = \sum_{j \ge k_{j_1+1}} \left( (UV)_{ij} + [f^{n_3}b, u_{ij}] \right) m_j
      + u \sum_j \left( (U \cdot \beta)_{ij} + (\alpha \cdot V)_{ij} \right) m_j
    \end{align*}

    \noindent Therefore,
    \[
      [f^{n_2}a, f^{n_3}b] m_i = \sum_{j \ge k_{j_1+1}} c_{ij} m_j + u \sum_j \gamma_{ij} m_j
    \]
    where
    \[
      c_{ij} = [V, U]_{ij} + [f^{n_2}a, v_{ij}] - [f^{n_3}b, u_{ij}],
    \]
    and
    \[
      [\gamma_{ij}] = [V, \alpha] + [\beta, U].
    \]
    \noindent On the other hand, since $\overline{T}$ is commutative
    \[
      [f^{n_2}a, f^{n_3}b] = u \cdot c \quad \text{for some } c \in T
    \]

    \noindent This implies
    \[
      [f^{n_2}a, f^{n_3}b] m_i = u \cdot c m_i = u \cdot \sum_j d'_{ij} m_j, \quad \overline{d'_{ij}} \in S_{\bar f}
    \]
    \noindent (Here we use the $E$-flatness of $\mathcal{M}$, that is the isomorphism $\overline{\underline{\mathcal{M}}} \xrightarrow{\sim} u \underline{\mathcal{M}}$.)

    \noindent Canceling the denominators of $d'_{ij}$, we set
    \[
      f^{n_1} [f^{n_2}a, f^{n_3}b] m_i = u \cdot \sum_j e_{ij} m_j, \quad \overline{e_{ij}} \in S
    \]
    So
    \[
      f^{n_1} \cdot [c_{ij}] + u f^{n_1} [\gamma_{ij}] = u [e_{ij}]
    \]
    But $[c_{ij}]$ is traceless because it is a sum of a commutator and a strict upper triangular matrix.
    $[\gamma_{ij}]$ is also traceless similarly.
    So $[\overline{e_{ij}}]$ is traceless.
  \end{proof}

  \noindent Now we go back to prove $\{\frak{p}, \frak{p}\} \subseteq \frak{p}$.\\

  \noindent We know
  $  f^{n_1} [f^{n_2}a, f^{n_3}b] = c $
  \begin{align*}
    c & = f^{n_1} f^{n_2} a f^{n_3} b - f^{n_1} f^{n_3} b f^{n_2} a                                          \\
      & = f^{n_1+n_2+n_3} ab + f^{n_1+n_2} [a, f^{n_3}] b  - f^{n_1+n_2+n_3} ba - f^{n_1+n_3} [b, f^{n_2}] a \\
      & = f^{n_1+n_2+n_3} [a, b] + f^{n_1+n_2} [a, f^{n_3}] b - f^{n_1+n_3} [b, f^{n_2}] a
  \end{align*}

  \noindent This implies $\overline{c} \in \bar{f}^{n_1+n_2+n_3} \{\bar{a}, \bar{b}\} + \frak{p}$.\\

  \noindent By Lemma~\ref{lem:matrix}, $\text{the trace of } \overline{c} \text{ on } \overline{\underline{\mathcal{M}}}_{\bar f} = \sum_i \overline{e_{ii}} = 0$.
  Since $\frak{p}^l_{\bar f} \cdot \overline{\underline{\mathcal{M}}}_{\bar f} = 0$, its trace is also 0.
  So the trace of $\{\bar{a}, \bar{b}\}$ on $\overline{\underline{\mathcal{M}}}_{\bar f}$ is also 0 (since $\bar{f} \notin \frak{p}$).\\

  \noindent Then $\forall x \in \overline{T}_{\bar f}$, $\{x\bar{a}, \bar{b}\} = x\{\bar{a}, \bar{b}\} + \{x, \bar{b}\}\bar{a} \in x\{\bar{a}, \bar{b}\} + \frak{p}_{\bar f}$.
  Similarly, we conclude
  \[
    \text{the trace of } x\{\bar{a}, \bar{b}\} \text{ on } \overline{\underline{\mathcal{M}}}_{\bar f} \text{ is also } 0.
  \]

  \noindent Meanwhile, since $\overline{\underline{\mathcal{M}}}_{\bar f}$ is free over $B = \overline{T}_{\bar f} / \frak{p}_{\bar f}$ and $S_{\bar f} \hookrightarrow B$ is a finite free extension, the trace of $y = x\{\bar{a}, \bar{b}\} \in \overline{T}_{\bar f}$ on $\overline{\underline{\mathcal{M}}}_{\bar f}$ is
  \[
    r \cdot (\text{the trace of } \bar{y} \text{ on } B),
  \]
  where $r$ is the rank of $\overline{\underline{\mathcal{M}}}_{\bar f}$ over $S_{\bar f}$.
  (For this properties of trace one can refer \cite[Theorem 7.8.5]{li2019algebraic})\\

  \noindent Since finite extension is integral and we are over $\mathbb{C}$ (thus the corresponding field extension is separable), the trace form of $S_{\bar f} \hookrightarrow B$ is non-degenerate.

  \noindent So
  \[
    \bar{y} = 0 \in B \implies \{\bar{a}, \bar{b}\} \in \frak{p}.
  \]

  This proves $\{\mathfrak p,\mathfrak p\}\subseteq \mathfrak p$, so $\mathfrak p$ (and thus $I$) is involutive.  The proof of
  Theorem~\ref{thm:gabber-II}---and hence Theorem~\ref{thm:gabber-I}---is complete.

\end{proof}